% !TeX root = ../深度学习笔记.tex
% 由主文档通过 \input 载入；请编译主文档。
\setcounter{chapter}{15} % 保留原书章号。
\chapter{迁移学习、预训练与微调}
\label{chap:dl-transfer}
从已学参数开始训练，可以利用源任务中形成的表示，也可能带来不适合目标任务的偏置。迁移是否有益，必须相对于相同目标数据、评价协议和适配预算下的基线判断。

\section{源域、目标域与迁移条件}
\emph{迁移学习（transfer learning）}利用源问题获得的信息帮助目标问题；\emph{预训练（pretraining）}是获得初始表示或参数的训练阶段，\emph{微调（fine-tuning）}是以这些参数为起点适配目标任务。源、目标可以具有不同输入分布、标签空间和训练损失，三种差异应分别说明。

例如源任务学习辨别物体形状，目标任务也依赖边缘和纹理时，部分表示可能有用；若源任务鼓励忽略颜色而目标恰要按颜色分类，这种不变性可能妨碍目标学习。相对合适基线表现下降称\emph{负迁移（negative transfer）}。能载入已有权重是程序事实，能够节省数据或提高准确率是需要实验支持的结论。

\section{冻结特征、部分微调与全量微调}
冻结特征时只训练新输出头，优化变量最少，适合检验源表示是否已经可用；部分微调只开放若干模块；全量微调更新全部可训练权重，表达调整能力更大，但小数据下也更容易过拟合。不同模块可使用不同学习率，选择应依据验证结果，而不是默认底层一定通用、顶层一定专用。

“冻结权重”与“推断模式”不同：前者禁止梯度更新，后者控制 dropout、归一化统计等行为。即使全部权重冻结，训练模式下的运行均值仍可能变化。实现时应明确可训练参数集合、缓冲状态和优化器状态，检查新输出头没有被遗漏，冻结模块也没有被意外更新。

\section{参数高效适配与低秩更新}
\emph{参数高效微调（parameter-efficient fine-tuning, PEFT）}仅训练相对少量新增或选定参数。适配器（adapter）可在冻结层旁加入小瓶颈分支；\emph{低秩适配（low-rank adaptation, LoRA）}则按文献\cite{hu2021}把冻结矩阵 $\mathbf W_0\in\mathbb R^{m\times d}$ 改为
\begin{equation}
\mathbf W=\mathbf W_0+\frac\alpha r\mathbf B\mathbf A,
\qquad \mathbf B\in\mathbb R^{m\times r},\quad\mathbf A\in\mathbb R^{r\times d}.
\label{eq:dl-low-rank}
\end{equation}
其中 $r\geq1$、$\alpha>0$，更新秩不超过 $r$，可训练参数为 $r(m+d)$。对输入 $\mathbf x$，先算 $\mathbf A\mathbf x$ 再乘 $\mathbf B$，不必在每次前向显式构造稠密增量。

令 $c=\alpha/r$，把损失对 $\mathbf W$ 的梯度记为 $\mathbf G$，则 $\nabla_{\mathbf B}\ell=c\mathbf G\mathbf A^\mathsf T$、$\nabla_{\mathbf A}\ell=c\mathbf B^\mathsf T\mathbf G$。若两因子同时初始化为零，两者梯度都为零；常用一因子随机、另一因子为零，使初始增量为零而保留学习通路。参数减少并不同比减少激活或全部前向计算，模型仍要通过冻结主干。

\section{分布变化、遗忘与选择协议}
目标微调后原任务能力下降称为\emph{灾难性遗忘（catastrophic forgetting）}的一种表现，但必须固定源任务评价才能识别。目标数据与源数据的类别比例变化、标签定义变化、条件关系变化是不同问题；只看目标准确率无法推断源能力保留程度。

比较从头训练、冻结表示和微调时，固定目标划分、标注量、主要选择预算与指标，分别报告使用的源模型和预训练信息。目标训练成本可以单独比较，总成本则还应说明是否计入预训练。少量标签场景可用多次独立划分考察变动，不能用测试集选择秩、解冻范围或训练轮次。

\section{例题与习题}
\begin{example}[低秩适配的参数量]
对 $512\times512$ 权重矩阵，取 LoRA 秩 $r=8$。比较全量微调和低秩适配的可训练参数量，不计偏置。
\end{example}
\begin{solve}
全量需 $512^2=262144$ 个参数，低秩需 $8(512+512)=8192$ 个，是前者的 $1/32$。原矩阵仍需存储和参与计算，故该比率只描述可训练参数，不能直接当作整体内存或延迟比率。
\end{solve}
\begin{enumerate}
\item 推导两个低秩因子的梯度，解释零初始化单个因子与同时零初始化的区别。
\item 对一个源、目标标签语义不同的任务，说明应保留哪些模块、替换哪些输出，并给出可验证假设。
\item 固定目标标签预算比较四种适配方式，同时报告目标指标、源任务指标、可训练参数和训练时间。
\end{enumerate}

